\documentclass[a4paper,12pt]{report}
\usepackage[margin=0.85in]{geometry}
\usepackage[utf8]{inputenc}
\usepackage{graphicx}
\usepackage{mathptmx}
\usepackage{parskip}
\usepackage{listings}
\usepackage{titlesec}
\usepackage{setspace}
\usepackage{float}
\usepackage{xcolor}
\usepackage{amssymb}
\usepackage{etoolbox}
\usepackage{fancyhdr}

\setlength{\parindent}{1cm}
\setlength{\parskip}{6pt}
\setstretch{1.5}
\setcounter{tocdepth}{3}
\setcounter{secnumdepth}{3}
\frenchspacing

% ---------- Student details (edit here only) ----------
\newcommand{\studentname}{BRIJESH SEBASTIAN}
\newcommand{\rollno}{CHN25MCA-2023}
\newcommand{\guidename}{Smt. Jithy John}

% ---------- Colours and code style (for a future code appendix) ----------
\definecolor{codegreen}{rgb}{0,0.6,0}
\definecolor{codegray}{rgb}{0.5,0.5,0.5}
\definecolor{codepurple}{rgb}{0.58,0,0.82}

\lstdefinestyle{customcsharp}{
	language=[Sharp]C,
	basicstyle=\small\ttfamily,
	keywordstyle=\color{blue},
	commentstyle=\color{codegreen},
	stringstyle=\color{orange},
	showstringspaces=false,
	breaklines=true,
	frame=single,
	numbers=left,
	numberstyle=\tiny\color{codegray},
	numbersep=5pt,
	tabsize=2,
	captionpos=b
}

% ---------- Headings ----------
\renewcommand{\chaptername}{CHAPTER}
\let\mtcontentsname\contentsname
\renewcommand\contentsname{\MakeUppercase\mtcontentsname}
\renewcommand{\listfigurename}{LIST OF FIGURES}
\renewcommand{\listtablename}{LIST OF TABLES}

\titleformat{\chapter}[display]
{\normalfont\sffamily\Large\bfseries\centering}
{\chaptertitlename\ \thechapter}{0pt}{\Large}

\titlespacing*{\chapter}{0pt}{0pt}{6pt}
\titlespacing*{\section}{0pt}{0pt}{6pt}
\titlespacing*{\subsection}{0pt}{0pt}{6pt}
\titlespacing*{\subsubsection}{0pt}{0pt}{6pt}

\titleformat*{\section}{\normalfont\sffamily\Large\bfseries}
\titleformat*{\subsection}{\normalfont\sffamily\large\bfseries}
\titleformat*{\paragraph}{\normalfont\large\sffamily}
\titleformat*{\subsubsection}{\normalfont\sffamily\large\itshape}

% ---------- Stop \cite from printing inside the lists of figures/tables ----------
\makeatletter
\let\oldcite\cite
\pretocmd{\listoffigures}{\def\cite{\ignorespaces\@gobble}}{}{}
\apptocmd{\listoffigures}{\let\cite\oldcite}{}{}
\pretocmd{\listoftables}{\def\cite{\ignorespaces\@gobble}}{}{}
\apptocmd{\listoftables}{\let\cite\oldcite}{}{}
\makeatother

% =====================================================================
%  IMAGE FILES USED (keep these exact names in the same folder as Main.tex)
%  CEC.png
%  usecase-saferoute.png        dfd-levle-0.png
%  dfd-citizen-dfd-levle-1.png  dfd-admin-dfd-level-1.png
%  dfd-responder-dfd-levle-1.png
%  system-arch.png      dbschema.jpeg
%  Screenshot-*.png  (11 screenshots, see the Screenshots section of Chapter 7)
% =====================================================================

\begin{document}
	
	% =====================================================================
	%  COVER PAGE
	% =====================================================================
	\begin{minipage}{0.9\linewidth}
		\centering
		
		\vspace{2cm}
		\Large\textbf{\uppercase{SAFEROUTE}}\\
		\vspace{0.3cm}
		\large\textbf{Emergency Response \& Incident Reporting Platform}\\
		\vspace{1cm}
		
		{\uppercase{\large{{MINI PROJECT REPORT}\par}}}
		\vspace{0.5cm}
		\large{Submitted by}
		
		\vspace{0.5cm}
		{\large \textbf{\studentname \\ \rollno}\par}
		\vspace{0.5cm}
		to
		\vspace{0.5cm}
		
		\textit{APJ Abdul Kalam Technological University \\ in partial fulfillment of the requirements for the award of Degree
			in \textbf{Master of Computer Applications}}
		
		\vspace{2cm}
		\includegraphics[width=0.3\textwidth]{CEC.png}
		\vspace{1cm}
		
		\large{\textbf{DEPARTMENT OF COMPUTER ENGINEERING\\
				COLLEGE OF ENGINEERING CHENGANNUR, ALAPPUZHA\\
				OCTOBER 2026\\}}
	\end{minipage}
	\thispagestyle{empty}
	\clearpage
	
	% =====================================================================
	%  CERTIFICATE
	% =====================================================================
	\thispagestyle{empty}
	
	\begin{center}
		\vspace{1.5cm}
		
		\textbf{DEPARTMENT OF COMPUTER ENGINEERING \\}
		\textbf{COLLEGE OF ENGINEERING CHENGANNUR\\}
		\textbf{ALAPPUZHA}
	\end{center}
	
	\begin{center}
		\includegraphics[width=0.3\textwidth]{CEC.png}
	\end{center}
	
	\begin{center}
		\textbf{CERTIFICATE}
	\end{center}
	
	\begin{spacing}{1.5}
		\textit{This is to certify that the project report titled \textbf{SAFEROUTE: EMERGENCY RESPONSE \& INCIDENT REPORTING PLATFORM} is a bonafide record of the \textbf{20MCA245 MINI PROJECT} presented by \textbf{\studentname} (\rollno), Third Semester Master of Computer Applications student, under our guidance and supervision. This project is submitted in partial fulfillment of the requirements for the award of the degree \textbf{Master of Computer Applications} of \textbf{APJ Abdul Kalam Technological University}.}
	\end{spacing}
	
	% Signature block
	\begin{singlespace}
		\vspace*{1cm}
		
		\noindent
		\begin{tabular*}{\textwidth}{@{\extracolsep{\fill}} l l}
			\textbf{\guidename} & \textbf{Smt. Josmi Jose} \\
			Project Guide & Project Coordinator \\
			Assistant Professor & Assistant Professor \\
			Dept. of Computer Engineering & Dept. of Computer Engineering \\[1.5cm]
		\end{tabular*}
		
		\begin{center}
			\textbf{Dr. Renu George} \\
			Professor \\
			Head of the Department \\
			Dept. of Computer Engineering
		\end{center}
	\end{singlespace}
	\clearpage
	
	\pagenumbering{roman}
	
	% =====================================================================
	%  DECLARATION
	% =====================================================================
	\addcontentsline{toc}{chapter}{Declaration}
	\begin{center}
		\Large {\bf \uppercase{Declaration}}
	\end{center}
	\vspace{1cm}
	
	\noindent I undersigned hereby declare that the project report ``\textbf{SAFEROUTE: EMERGENCY RESPONSE \& INCIDENT REPORTING PLATFORM}'', submitted for partial fulfillment of the requirements for the award of degree of Master of Computer Application of the APJ Abdul Kalam Technological University, Kerala is a bonafide work done by me under the supervision of \textbf{\guidename}, Assistant Professor, Department of Computer Engineering. This submission represents my ideas in my own words, and where ideas or words of others have been included, I have adequately and accurately cited and referenced the original sources. I also declare that I have adhered to the ethics of academic honesty and integrity and have not misrepresented or fabricated any data, idea, fact, or source in this submission. I understand that any violation of the above will be a cause for disciplinary action by the institute and/or the University and can also evoke penal action from the sources which have thus not been properly cited or from whom proper permission has not been obtained. This report has not been previously formed the basis for the award of any degree, diploma, or similar title of any other University.
	
	\vspace{4cm}
	
	\noindent\begin{minipage}{0.45\linewidth}
		\begin{flushleft}
			\textbf{Place:} Chengannur \\
			\textbf{Date :}
		\end{flushleft}
	\end{minipage}
	\hfill
	\begin{minipage}{0.45\linewidth}
		\begin{flushright}
			\textbf{\studentname}\\
			\textbf{\rollno}\\
		\end{flushright}
	\end{minipage}
	\clearpage
	
	% =====================================================================
	%  ACKNOWLEDGEMENT
	% =====================================================================
	\addcontentsline{toc}{chapter}{Acknowledgement}
	\begin{center}
		\Large {\bf \uppercase{Acknowledgement}}
	\end{center}
	\vspace{1.5cm}
	
	\noindent This work would not have been possible without the support of many people.
	First and foremost, I give thanks to Almighty God who gave me the inner strength, resources, and ability to complete my project successfully.\\
	I would like to thank \textbf{Dr. Hari V.S}, The Principal, College of Engineering Chengannur, for providing the best facilities and atmosphere for the project completion and presentation. Special thanks to \textbf{Dr. Renu George}, HOD, Department of Computer Engineering, for her exceptional support, guidance, and encouragement throughout the project. I would also like to thank my project coordinators \textbf{Smt. Josmi Jose}, Assistant Professor, Department of Computer Engineering, \textbf{Smt. Santhy Viswam}, Assistant Professor, Department of Computer Engineering, \textbf{Smt. Celin T}, Assistant Professor, Department of Computer Engineering, for their extended help and support during the project. I would also like to thank \textbf{\guidename}, Assistant Professor, Department of Computer Engineering, who also took on the role of my project guide, for her extended help and support during the project.\\
	I would like to thank my dear friends and faculty for extending their cooperation and encouragement throughout the project work, without which I would never have completed the project this well. Thank you all for your love and also for being very understanding.
	
	\begin{flushright}
		\textbf{\studentname\\\rollno}
	\end{flushright}
	\clearpage
	
	% =====================================================================
	%  ABSTRACT
	% =====================================================================
	\addcontentsline{toc}{chapter}{Abstract}
	\begin{center}
		\Large\textbf{ABSTRACT}
	\end{center}
	\vspace{1cm}
	
	\textbf{``SafeRoute''} is a web-based Emergency Response and Incident Reporting Platform developed to connect citizens, administrators, and emergency responders through a single centralized system. During emergencies, delays in reporting, unverified information, and poor coordination between the people who report an incident, the authority that dispatches help, and the teams that respond on the ground can increase the time taken to reach the affected location. SafeRoute addresses these challenges by providing role-based access for three types of users: reporters, administrators, and responders. Citizens can register, log in, and report emergencies by providing the incident type, severity, description, location, and supporting photographs. Administrators can review reported incidents, verify or reject them, and assign a suitable ambulance, police, or fire responder to each verified incident. Responders receive the assignment, view the incident location on a map, open it in Google Maps for navigation, accept or reject the dispatch, share their live location, and update the response status as they travel, arrive, and resolve the incident. Every change of status is recorded in an incident history, and in-app notifications keep reporters, administrators, and responders informed. The administrator module also provides dashboards, analytics, monthly reports, and user management. The backend is developed as a REST API using Node.js, Express.js, and Sequelize with a MySQL database, and it uses JWT-based authentication with refresh-token rotation, password hashing, request validation, rate limiting, and secured image delivery. The web interface is developed using React and Tailwind CSS, and maps are displayed using Leaflet with OpenStreetMap. By combining incident reporting, verification, dispatch, tracking, and history in one platform, SafeRoute aims to improve communication, coordination, and response tracking during emergencies.
	
	\clearpage
	
	% =====================================================================
	%  CONTENTS, LISTS
	% =====================================================================
	\tableofcontents
	\newpage
	
	\addcontentsline{toc}{chapter}{List of Figures}
	\listoffigures
	\newpage
	
	\addcontentsline{toc}{chapter}{List of Tables}
	\listoftables
	
	% =====================================================================
	%  CHAPTER 1
	% =====================================================================
	\pagenumbering{arabic}
	\chapter{INTRODUCTION}
	
	\section{Project Area}
	\par
	Emergency response depends on how quickly accurate information reaches the people who can act on it. When an accident, fire, medical emergency, or other hazard occurs, the first information usually comes from a citizen at the scene. That information must be checked, passed to a suitable emergency unit, and followed until the incident is resolved. In many places this chain still relies on phone calls, messaging applications, and informal communication, which makes it difficult to verify reports, to know which responder has been sent, and to track what is happening after a report is made.
	
	``\textbf{SafeRoute}'' is proposed as a web-based Emergency Response and Incident Reporting Platform that provides a centralized system for reporting, managing, assigning, and tracking emergency incidents. Citizens can report incidents with important details such as the incident type, severity, description, location, and photographs. Administrators can verify incidents, prioritize them, and assign suitable responders. Responders can receive assignments, view the incident location, navigate to it, and update their response status until the incident is resolved. The system maintains a centralized database of incidents, users, assignments, and status updates, so that every action taken on an incident is recorded.
	
	Developed using web technologies, a Node.js and Express.js REST API, and a MySQL database, SafeRoute combines role-based access control, location services, notifications, and incident history in one platform. It aims to improve communication, coordination, response tracking, and incident management during emergencies.
	
	\section{Objectives}
	\begin{enumerate}
		\item \textbf{To develop a centralized platform for emergency incident management:}
		The system aims to provide a single web-based platform where emergency incidents can be reported, verified, assigned, tracked, and resolved, instead of depending on scattered communication channels.
		
		\item \textbf{To enable quick and detailed incident reporting:}
		The system aims to allow citizens to report emergencies quickly with essential details such as the incident type, severity, description, location on a map, and photographs that serve as supporting evidence.
		
		\item \textbf{To support verification and prioritization of incidents:}
		Administrators should be able to review every reported incident, verify genuine reports, reject false or duplicate reports, and use the severity of an incident to decide how urgently a response is needed.
		
		\item \textbf{To assign suitable responders to verified incidents:}
		The system aims to let administrators select an available ambulance, police, or fire responder for each verified incident, and to reassign or cancel the assignment when required.
		
		\item \textbf{To support responders during the response:}
		Responders should be able to view the assigned incident, see its location, open it in a navigation application, accept or reject the assignment, share their live location, and update their status when they start travelling, arrive at the location, and resolve the incident.
		
		\item \textbf{To provide real-time status tracking and notifications:}
		Reporters, administrators, and responders should be able to follow an incident from reporting to resolution, and the system should notify the concerned users whenever an important status change takes place.
		
		\item \textbf{To maintain a complete history of incidents and response activities:}
		The system aims to record every status change together with the user who made it and the time at which it happened, so that incident records are available for monitoring, reporting, and analysis.
		
		\item \textbf{To provide secure, role-based access:}
		The system aims to ensure that reporters, administrators, and responders can access only the functions meant for their roles, and that user data and incident data are protected.
	\end{enumerate}
	
	% =====================================================================
	%  CHAPTER 2
	% =====================================================================
	\chapter{LITERATURE REVIEW}
	
	\section{Advancing Smart Disaster Response by Leveraging Social Sensing and Mobile Technology}
	
	\hspace{1cm}This paper focuses on using mobile technology and social sensing to improve the collection of information during disasters. The authors developed an agile, iterative, and user-centred mobile Geographic Information System (GIS) application and tested it during real hurricane events. The application enabled real-time and offline reporting with GPS location, categorized information, and multimedia, so that people affected by a disaster could contribute information from the field. The main limitation of the work is that it concentrates on collecting and visualizing disaster information rather than on complete responder coordination. This paper is relevant to SafeRoute because it supports the idea of collecting incident reports together with GPS location, categories, and photographs. SafeRoute extends this idea by adding the steps that follow reporting, namely verification by an administrator, assignment of a responder, and tracking of the response until resolution.~\cite{one}
	
	\vspace{0.5cm}
	\section{From Dispatch to Impact: Evaluating ICT-enabled Co-production of Emergency Response}
	
	\hspace{1cm}This paper studies how mobile applications can be used to dispatch volunteer first responders to emergencies. The authors used a mixed-method approach to study volunteer first responders who were dispatched through mobile applications. The results show that mobile dispatch helped volunteers arrive before professional responders, which reduced the first-response time and improved incident management. The paper also notes that the approach depends on trained volunteers, reliable information and communication technology (ICT), equipment, and coordination, and that delayed alerts and application failures can affect performance. This paper is relevant to SafeRoute because it shows the value of dispatching responders through a mobile or web application and highlights the importance of reliable notifications. In SafeRoute, assignments are issued by an administrator after verification and are delivered to responders as in-app notifications and assignment lists.~\cite{two}
	
	\vspace{0.5cm}
	\section{`Invented-on-the-fly' Mobile Application for Disaster Response: Construction of Technological Frames and Impact}
	
	\hspace{1cm}This paper is a case study that examines how a mobile disaster-response application was developed and adapted during the Gaja cyclone. The study shows that disaster-response technology needs to adapt to changing requirements, users, and emergency workflows. Its limitation is that it is based on a specific disaster and organizational context, so its findings may not directly generalize to every emergency-response environment. This paper is relevant to SafeRoute because it shows that emergency applications must fit the way different users actually work during a response. SafeRoute therefore separates the workflows of reporters, administrators, and responders, and manages them through clear incident and assignment statuses that can be extended in the future.~\cite{three}
	
	\vspace{0.5cm}
	\section{Engineering for Crisis Management: A User-Centred Analysis of Disaster Mobile Applications}
	
	\hspace{1cm}This paper analyses existing disaster mobile applications from the user's point of view. The authors analysed 70 disaster mobile apps and user reviews using feature analysis, topic modelling, and sentiment analysis. They found major issues involving reliability, usability, accessibility, and information clarity, and observed that most apps focus more on response than on preparation and recovery. The study concentrates mainly on the analysis of existing disaster apps, rather than on developing an integrated reporting, verification, and dispatch system. This paper is relevant to SafeRoute because it identifies the quality factors that an emergency application must satisfy. SafeRoute is designed with a simple interface, clear status information, and a responsive layout so that reporters and responders can use it with minimal steps.~\cite{four}
	
	\vspace{0.5cm}
	\section{Gap Identified}
	
	\hspace{1cm}The reviewed work shows that existing systems focus mainly on incident reporting and information collection, and that many solutions do not provide a complete end-to-end emergency response workflow. Reporter, administrator, and responder activities are often handled separately, and there is limited integration between incident verification, responder assignment, and status tracking. Existing applications also show issues with usability, reliability, accessibility, and information clarity. There is therefore a need for a centralized platform that connects reporting, coordination, dispatch, and resolution. SafeRoute addresses this gap through a role-based, centralized emergency incident management system.
	
	\clearpage
	
	% =====================================================================
	%  CHAPTER 3
	% =====================================================================
	\chapter{PROBLEM DEFINITION AND MOTIVATIONS}
	
	\section{Existing System}
	
	In the existing approach, emergencies are usually reported by phone calls, messaging applications, social media, or through personal contacts, and the response is coordinated manually. Reports reach the authorities in different formats and with different levels of detail, and it is difficult to verify them or to know whether a responder has already been sent. The reporter often receives no information about the progress of the response, while responders may not receive the exact location or complete details of the incident. Existing digital solutions focus mainly on collecting and displaying incident information and do not provide a complete workflow that connects reporting, verification, assignment, response, and resolution. SafeRoute addresses these limitations by providing a centralized web-based platform where incidents are reported with their location and evidence, verified by an administrator, assigned to a suitable responder, and tracked until they are resolved.
	
	\section{Limitations}
	\begin{itemize}
		\item \textbf{No Centralized Platform :} Incident information is spread across calls, messages, and separate systems, which makes coordination difficult.
		
		\item \textbf{Incomplete Incident Information :} Informal reports may not contain the exact location, the type and severity of the incident, or photographic evidence.
		
		\item \textbf{No Verification Step :} Unverified and duplicate reports can lead to wasted effort and delays in dispatching help.
		
		\item \textbf{Delayed Dispatch :} Manual assignment of responders and passing of information from person to person delay the response.
		
		\item \textbf{Separate Handling of Roles :} Reporter, administrator, and responder activities are often handled separately, with limited integration between them.
		
		\item \textbf{Limited Status Tracking :} Reporters usually cannot see whether an incident has been verified, assigned, or resolved.
		
		\item \textbf{Difficulty in Reaching the Location :} Responders may not receive the incident location in a form that can be used directly for navigation.
		
		\item \textbf{Lack of Notifications :} Users are not informed automatically when an assignment is made or when the status of an incident changes.
		
		\item \textbf{No Incident History :} Records of who did what and when are often missing, which limits monitoring, transparency, and analysis.
		
		\item \textbf{Usability and Reliability Issues :} Existing applications show issues with usability, reliability, accessibility, and information clarity.
		
		\item \textbf{Limited Administrative Control :} Informal methods provide no dedicated facilities for managing users, responders, and incident records.
	\end{itemize}
	
	\section{Problem Statement}
	\par To eliminate emergency response delays and coordination gaps through a centralized platform for real-time reporting, tracking, and responder management.
	
	\section{Proposed System}
	\par SafeRoute is a web-based Emergency Response and Incident Reporting Platform that connects reporters, administrators, and responders through a single system. Reporters submit emergencies with their type, severity, description, location, and photographs. Administrators verify the reports, prioritize them, and assign suitable responders. Responders accept the assignment, navigate to the incident location, share their live location, and update the status until the incident is resolved. The system sends notifications for important events, maintains a history of every status change, and provides the administrator with dashboards, analytics, and reports. Role-based access control ensures that every user can use only the functions that belong to their role.
	
	The key features of the proposed system include:
	\begin{enumerate}
		\item \textbf{User Registration and Login :} Allows citizens to create accounts and securely log in, with separate access for reporters, administrators, and responders.
		
		\item \textbf{Incident Reporting :} Allows reporters to submit emergencies with the incident type, severity, description, location, and up to five photographs.
		
		\item \textbf{Location Selection :} Detects the reporter's current location and allows it to be changed by searching for a place or moving a pin on the map.
		
		\item \textbf{Incident Verification and Prioritization :} Allows administrators to verify or reject reported incidents and to use severity to prioritize them.
		
		\item \textbf{Responder Assignment :} Allows administrators to assign an ambulance, police, or fire responder to a verified incident, and to reassign or cancel the assignment.
		
		\item \textbf{Response Management :} Allows responders to accept or reject assignments and to update the status as they start travelling, arrive, and resolve the incident.
		
		\item \textbf{Navigation Support :} Provides a button that opens the incident location in Google Maps so that responders can reach the location easily.
		
		\item \textbf{Live Responder Location :} Allows the responder's location to be shared so that reporters and administrators can follow the response on a map.
		
		\item \textbf{Notifications :} Informs users about assignments and status changes through in-app notifications.
		
		\item \textbf{Incident History and Reports :} Maintains the history of every status change and provides dashboards, analytics, and monthly reports for administrators.
		
		\item \textbf{User Management :} Allows administrators to create responder accounts and to view, edit, and disable users.
	\end{enumerate}
	
	% =====================================================================
	%  CHAPTER 4
	% =====================================================================
	\chapter{REQUIREMENT ANALYSIS}
	
	\section{Hardware Requirements}
	\begin{itemize}
		\item \textbf{Processor:} Intel Core i3 / equivalent or higher
		\item \textbf{Memory (RAM):} Minimum 4 GB
		\item \textbf{Storage:} Minimum 10 GB free space
		\item \textbf{Display:} 1366 $\times$ 768 resolution or higher
		\item \textbf{Network:} Wi-Fi or Ethernet connection
		\item \textbf{Input Devices:} Keyboard and mouse
		\item \textbf{Mobile Device:} Android/iOS smartphone for mobile access and location services
		\item \textbf{GPS:} Required on responder and reporting devices for location-based features
	\end{itemize}
	
	\section{Software Requirements}
	\begin{itemize}
		
		\item \textbf{Operating System:} Windows 10 / Windows 11, Linux, or macOS
		
		\item \textbf{Programming Language:}
		\begin{itemize}
			\item \textbf{JavaScript:} Used for both backend development (Node.js) and frontend development (React)
			\item \textbf{SQL:} Used for database definition and queries
		\end{itemize}
		
		\item \textbf{Framework:}
		\begin{itemize}
			\item \textbf{Express.js:} Used to build the REST API and handle routing, middleware, requests, and responses
			\item \textbf{React:} Used to build the single-page web interface
			\item \textbf{Tailwind CSS:} Frontend framework used to create a responsive user interface
		\end{itemize}
		
		\item \textbf{Database:}
		\begin{itemize}
			\item \textbf{MySQL:} Used to store and manage user, incident, assignment, notification, and history data.
		\end{itemize}
		
		\item \textbf{Libraries:}
		\begin{itemize}
			\item \textbf{Sequelize and mysql2:} Used for object-relational mapping and for connecting the application to MySQL.
			\item \textbf{jsonwebtoken and bcrypt:} Used for token-based authentication and password hashing.
			\item \textbf{Zod:} Used to validate request data.
			\item \textbf{Multer:} Used to handle photograph and profile image uploads.
			\item \textbf{Helmet, CORS, and express-rate-limit:} Used for security headers, cross-origin access control, and request rate limiting.
			\item \textbf{Nodemailer:} Used to send password reset links by email.
			\item \textbf{React Router, TanStack Query, and Axios:} Used for page navigation, data fetching, and communication with the REST API.
			\item \textbf{Leaflet and OpenStreetMap:} Used to display maps and incident locations.
			\item \textbf{Recharts:} Used to display analytics charts for administrators.
		\end{itemize}
		
		\item \textbf{Development Tools:}
		\begin{itemize}
			\item \textbf{Visual Studio Code:} For coding and development
			\item \textbf{Postman:} Used for testing the REST API
			\item \textbf{Git:} Used for version control
			\item \textbf{Web Browser (Chrome/Edge/Firefox):} For running and testing the application
		\end{itemize}
		
		\item \textbf{Internet Connection:} Required for location services, notifications, and online system access.
		
	\end{itemize}
	
	\section{System Features}
	\begin{enumerate}
		
		\item \textbf{User Registration and Login:}
		Allows users to create accounts and log in securely. Reporters register themselves, responder accounts are created by an administrator, and the administrator account is created during system setup. Passwords are stored as hashes.
		
		\item \textbf{Role-Based Access:}
		Provides separate functions and screens for reporters, administrators, and responders, and prevents users from accessing functions that do not belong to their role.
		
		\item \textbf{Incident Reporting:}
		Allows reporters to submit new emergency incidents by providing the incident type, severity, description, location, and supporting evidence.
		
		\item \textbf{Evidence Upload:}
		Allows reporters to attach up to five photographs (JPEG, PNG, or WebP) to an incident. Uploaded files are checked and are available only to authorized users.
		
		\item \textbf{Incident Verification and Management:}
		Allows administrators to view, search, filter, verify, and reject reported incidents and to prioritize them by severity.
		
		\item \textbf{Responder Assignment:}
		Allows administrators to assign a verified incident to a suitable available responder, and to reassign or cancel an assignment.
		
		\item \textbf{Response Management:}
		Allows responders to view their assigned incidents, accept or reject them, and update the status as they start travel, arrive, and resolve the incident.
		
		\item \textbf{Location Management and Navigation:}
		Records and displays incident locations on a map, shares the responder's location, and allows responders to open the incident location in Google Maps.
		
		\item \textbf{Incident Tracking and History:}
		Allows users to track the status of an incident, and maintains a history of every status change with the user and the time.
		
		\item \textbf{Notifications:}
		Sends in-app notifications for important incident and assignment updates, which can be viewed, marked as read, or deleted.
		
		\item \textbf{Dashboards, Analytics, and Reports:}
		Provides administrators with incident counters, analytics by type, severity, status, and month, a nearby-incident map, and monthly reports.
		
		\item \textbf{User Management:}
		Allows administrators to create responder accounts and to view, edit, and disable users.
		
		\item \textbf{Profile Management:}
		Allows users to update their details, upload a profile image, change their password, and log out of all devices.
		
		\item \textbf{Password Recovery:}
		Allows users who forget their password to receive a time-limited reset link by email.
		
		\item \textbf{Responsive User Interface:}
		Provides a user-friendly interface that adapts to desktops, laptops, tablets, and smartphones.
		
	\end{enumerate}
	
	\section{Non-Functional Requirements}
	
	\subsection{Performance Requirements}
	
	\begin{enumerate}
		\item \textbf{Fast Response:} The system should respond quickly to user requests.
		
		\item \textbf{Timely Updates:} Incident updates should be processed and shown without significant delay.
		
		\item \textbf{Efficient Database Access:} The system should retrieve and store data efficiently in the MySQL database.
		
		\item \textbf{Image Handling:} Photographs should be uploaded and displayed without significantly slowing down the system.
		
		\item \textbf{Controlled Polling:} Status and notification updates should be refreshed at suitable intervals so that the number of requests to the server remains within the configured limits.
		
		\item \textbf{Multiple Users:} The system should support multiple users working at the same time.
	\end{enumerate}
	
	\subsection{Quality Requirements}
	
	\begin{enumerate}
		\item \textbf{Usability:}
		The interface should be simple and easy to understand, and users should be able to report and manage incidents with minimal steps.
		
		\item \textbf{Reliability:}
		The system should operate consistently without frequent failures, and incident data should not be lost during operations.
		
		\item \textbf{Security:}
		User information and incident data must be protected, and only authorized users should access specific functions.
		
		\item \textbf{Availability:}
		The system should be accessible whenever emergency reporting is required.
		
		\item \textbf{Scalability:}
		The system should support increasing numbers of users and incidents.
		
		\item \textbf{Maintainability:}
		The system should be easy to maintain, update, and troubleshoot.
		
		\item \textbf{Compatibility:}
		The system should work across common browsers and supported devices.
		
		\item \textbf{Data Integrity:}
		Incident, user, assignment, and response information should remain accurate and consistent.
	\end{enumerate}
	
	% =====================================================================
	%  CHAPTER 5
	% =====================================================================
	\chapter{DESIGN AND IMPLEMENTATION}
	
	The proposed system is a web-based application named SafeRoute, designed to connect citizens, administrators, and emergency responders during emergencies. The system provides a centralized platform where incidents are reported with their location and evidence, verified by an administrator, assigned to a responder, and tracked until they are resolved. Role-based access, notifications, incident history, and dashboards make the response process more organized, transparent, and efficient.
	
	\section{Overall Design}
	SafeRoute is a web-based platform designed to report, verify, assign, and track emergency incidents. It uses React and Tailwind CSS for the user interface, a Node.js and Express.js REST API for server-side processing, and MySQL for data storage. The system allows reporters to register and report incidents with photographs and location, administrators to verify incidents and assign responders, and responders to accept assignments, navigate to the incident, share their location, and update the response status. Maps are displayed using Leaflet with OpenStreetMap. Notifications and a status history keep all users informed about the progress of an incident. Overall, the system aims to provide a simple, secure, and organized platform that improves communication, coordination, and response tracking during emergencies.
	
	\section{System Design}
	The system is structured into three main layers: presentation layer, application layer, and data layer.
	The presentation layer is developed using React and Tailwind CSS, allowing users to register, log in, report incidents, view maps, manage assignments, and view dashboards according to their roles. The application layer uses a Node.js and Express.js REST API to handle user authentication, role-based authorization, incident management, responder assignment, notifications, and file handling. The data layer consists of a MySQL database that stores user, incident, assignment, history, and notification information, together with an upload directory for photographs.
	
	When a user interacts with the system, the request is sent to the REST API, which authenticates the user, checks the role, validates the data, processes the request, and returns a JSON response to the user interface. Each important action, such as verifying an incident or accepting an assignment, changes the status of the incident, records an entry in the status history, and creates notifications for the concerned users. This layered design ensures organized processing, maintainability, and efficient data management.
	
	\subsection{Client-Side Design}
	The client-side represents the user interface of the SafeRoute system. It is developed as a single-page application using React, with Tailwind CSS used to provide a simple, soft, and responsive design that adapts to desktops, tablets, and smartphones.
	
	The client-side includes pages such as login, registration, password recovery, incident reporting, my incidents, incident details and tracking, the responder dashboard and assignment pages, the administrator dashboard, incident management, nearby-incident map, analytics, reports, user management, notifications, and profile. Page navigation is managed by the router, and the interface changes according to the role of the logged-in user. The access token is kept in memory and is renewed automatically using a refresh token, so that users do not have to log in repeatedly.
	
	Maps are displayed using Leaflet with OpenStreetMap tiles. The reporting form detects the user's current location and also allows the location to be changed by typing the name of a place or by moving the pin on the map. Photographs are loaded through authenticated requests so that only authorized users can view them. Responders are provided with a button that opens the incident location in Google Maps for navigation.
	
	\subsection{Server-Side Design}
	The server-side is responsible for handling the application logic, processing user requests, and communicating with the database. It is developed as a REST API using Node.js and Express.js, with Sequelize as the object-relational mapper for MySQL.
	
	The server processes operations such as user registration and login, token refresh, incident creation and updates, image upload, incident verification, responder assignment, status updates, location updates, notifications, analytics, and user management. Requests pass through a pipeline of security headers, cross-origin and rate-limit checks, authentication, role authorization, and schema validation before they reach the business logic. State changes such as responder assignment are performed inside database transactions so that the data remains consistent. Errors are returned in a standard format, and events are recorded in log files.
	
	\section{System Architecture}
	The SafeRoute system follows a layered architecture consisting of Users, User Interface, Application Layer, and Data Layer.
	
	Users include Reporters, Administrators, and Responders, who interact with the system through a web-based user interface. A reporter submits an incident and views its status, an administrator manages incidents and assigns responders, and a responder receives assignments and updates the response status. The application layer implements the main functions of the system, which are user authentication and role management, incident management, assignment and tracking, notifications, and reports and history.
	
	The data layer uses a MySQL database to store users, incidents, assignments, status updates, notifications, and history. Photographs of incidents and profile images are stored in the application's upload directory and are served only to authorized users.
	
	\vspace{1.5cm}
	
	\begin{figure}[h!]
		\centering
		\includegraphics[width=1.0\linewidth,height=0.9\linewidth,keepaspectratio]{system-arch.png}
		\caption{System Architecture}
		\label{fig:architecture}
	\end{figure}
	\newpage
	
	\section{Use Case Diagram}
	\par
	The Use Case Diagram represents the interaction between the users and the SafeRoute system. It identifies the main users of the system and the activities they can perform. In this system, there are three primary actors: \textbf{Reporter}, \textbf{Admin}, and \textbf{Responder}.
	
	\textbf{Actors:}
	
	\begin{enumerate}
		\item \textbf{Reporter / Citizen:} The reporter is a citizen who can register and log in to the system, report an incident with its location and evidence, view the status of the reported incidents, and receive notifications.
		
		\item \textbf{Admin / Dispatcher:} The administrator is responsible for managing the system by managing users, verifying incidents, assigning responders, and viewing reports.
		
		\item \textbf{Responder:} The responder is an emergency unit member (ambulance, police, or fire) who can view the assigned incident, update the response status, and receive notifications.
	\end{enumerate}
	
	\begin{figure}[H]
		\centering
		\includegraphics[width=1.0\linewidth,height=0.55\textheight,keepaspectratio]{usecase-saferoute.png}
		\caption{Use Case Diagram for SafeRoute}
		\label{fig:usecase}
	\end{figure}
	
	\section{Data Flow Diagrams (DFD)}
	
	\par
	The Data Flow Diagram (DFD) illustrates how data flows within the SafeRoute system between reporters, administrators, responders, the application, and the database. Reporters send incident reports with their location and evidence to the system and receive report confirmation and status updates. Administrators send verification, priority, and assignment decisions and receive the incident queue, responder status, and updates. Responders receive incident assignments with the location and details, and send back their acceptance or rejection, status updates, and the resolution of the incident.
	
	The Level 0 diagram represents the whole system as a single process, the SafeRoute Emergency Response System, together with the three external entities. The Level 1 diagrams are drawn separately for the three users, the Reporter (Citizen), the Admin, and the Responder, and together they divide the system into five processes: user authentication and management, incident reporting and management, responder assignment and response management, schedule assignment, and reports and analytics. All processes read from and write to a central database, which stores user data, incident data, assignment details, response updates, and the data needed for reports.
	
	The DFD clearly shows how user requests are received and processed, how incidents move from reporters to administrators and then to responders, and how the status updates of the responders return to the reporters and administrators.
	
	\vspace{0.3cm}
	
	\begin{figure}[H]
		\centering
		\includegraphics[width=0.9\linewidth,height=0.5\textheight,keepaspectratio]{dfd-levle-0.png}
		\caption{DFD Level 0}
	\end{figure}
	
	\begin{figure}[H]
		\centering
		\includegraphics[width=0.9\linewidth,height=0.55\textheight,keepaspectratio]{dfd-citizen-dfd-levle-1.png}
		\caption{DFD Level 1 -- Reporter (Citizen)}
	\end{figure}
	
	\begin{figure}[H]
		\centering
		\includegraphics[width=0.9\linewidth,height=0.55\textheight,keepaspectratio]{dfd-admin-dfd-level-1.png}
		\caption{DFD Level 1 -- Admin}
	\end{figure}
	
	\begin{figure}[H]
		\centering
		\includegraphics[width=0.9\linewidth,height=0.55\textheight,keepaspectratio]{dfd-responder-dfd-levle-1.png}
		\caption{DFD Level 1 -- Responder}
	\end{figure}
	
	\clearpage
	
	\section{Database Design}
	
	\subsection{Database Requirements}
	The database of SafeRoute is designed to meet the following requirements:
	
	\begin{enumerate}
		\item \textbf{User Management:} Store user details, roles, login credentials, and account status, and support the reporter, administrator, and responder roles.
		
		\item \textbf{Incident Management:} Store the incident type, description, severity, location, status, and reported time, and maintain the reporter associated with each incident.
		
		\item \textbf{Responder Allocation:} Store the responder assignment for each incident, and track the assignment status, assigned time, and completion time.
		
		\item \textbf{Incident Updates:} Store status changes and response updates made during an incident, and maintain the history of actions performed on an incident.
		
		\item \textbf{Media Management:} Store references to the images attached to incidents.
		
		\item \textbf{Notification Management:} Store notifications sent to users regarding incidents, assignments, and status changes.
		
		\item \textbf{Data Integrity and Security:} Use unique identifiers and relationships between records, protect user data, restrict access according to user roles, and maintain consistent and reliable incident records.
	\end{enumerate}
	
	\subsection{Database Schema}
	The relational schema consists of tables for users, incidents, incident images, incident assignments, incident status history, notifications, refresh tokens, and password reset tokens. A reporter can report many incidents, an incident can have many images, assignments, and status history entries, and a user can receive many notifications. Table~\ref{tab:dbtables} summarizes the tables, and Figure~\ref{fig:schema} shows the relationships between them.
	
	\begin{table}[h!]
		\centering
		\resizebox{\textwidth}{!}{
			\begin{tabular}{|l|p{6.5cm}|p{7.5cm}|}
				\hline
				\textbf{Table} & \textbf{Purpose} & \textbf{Key Attributes} \\
				\hline \hline
				\texttt{users} & Stores reporters, responders, and administrators & id, full\_name, email, phone, password (hash), role, responder\_type, profile\_image, is\_active, last\_known\_latitude, last\_known\_longitude \\
				\hline
				\texttt{incidents} & Stores the reported emergencies & id, title, description, incident\_type, severity, status, latitude, longitude, address, reporter\_id \\
				\hline
				\texttt{incident\_images} & Stores references to photographs attached to incidents & id, incident\_id, image\_url \\
				\hline
				\texttt{incident\_assignments} & Stores the dispatch of a responder to an incident & id, incident\_id, responder\_id, assigned\_by, assignment\_status, notes, responded\_at, completed\_at \\
				\hline
				\texttt{incident\_status\_history} & Stores every status change of an incident & id, incident\_id, previous\_status, new\_status, changed\_by, notes \\
				\hline
				\texttt{notifications} & Stores in-app notifications for users & id, user\_id, title, message, type, reference\_id, is\_read \\
				\hline
				\texttt{refresh\_tokens} & Stores hashed refresh tokens for login sessions & id, user\_id, token (hash), expires\_at \\
				\hline
				\texttt{password\_reset\_tokens} & Stores hashed, single-use password reset tokens & id, user\_id, token\_hash, expires\_at \\
				\hline
		\end{tabular}}
		\caption{Database Tables}
		\label{tab:dbtables}
	\end{table}
	
	\begin{figure}[H]
		\centering
		\includegraphics[width=1.0\linewidth,height=0.55\textheight,keepaspectratio]{dbschema.jpeg}
		\caption{Database Schema}
		\label{fig:schema}
	\end{figure}
	
	\subsection{Incident Lifecycle}
	Every incident follows a defined lifecycle. The status of an incident can be changed only by the user and in the order shown in Table~\ref{tab:lifecycle}, and every change is recorded in the status history.
	
	\begin{table}[h!]
		\centering
		\resizebox{\textwidth}{!}{
			\begin{tabular}{|l|l|l|l|}
				\hline
				\textbf{Current Status} & \textbf{Action} & \textbf{Performed By} & \textbf{New Status} \\
				\hline \hline
				Pending & Verify incident & Admin & Verified \\
				\hline
				Pending & Reject incident & Admin & Rejected \\
				\hline
				Pending / Verified & Cancel incident & Reporter & Cancelled \\
				\hline
				Verified & Assign responder & Admin & Assigned \\
				\hline
				Assigned & Accept assignment & Responder & Accepted \\
				\hline
				Assigned & Reject assignment & Responder & Verified \\
				\hline
				Assigned / Accepted & Reassign or cancel assignment & Admin & Verified \\
				\hline
				Accepted & Start travel & Responder & On the way \\
				\hline
				On the way & Arrived at location & Responder & Reached \\
				\hline
				Reached & Resolve incident & Responder & Resolved \\
				\hline
		\end{tabular}}
		\caption{Incident Status Transitions}
		\label{tab:lifecycle}
	\end{table}
	
	\clearpage
	
	\section{Module Description}
	
	\noindent \textbf{SafeRoute} is designed as a role-based web application that enables citizens to report emergencies, administrators to verify and dispatch, and responders to respond. It operates through three primary modules, \textbf{Reporter Module}, \textbf{Admin Module}, and \textbf{Responder Module}, which work together to manage an incident from reporting to resolution.
	
	\subsection{Reporter Module}
	The Reporter Module allows citizens to report emergencies and follow their progress.
	
	\begin{enumerate}
		\item \textbf{Registration and Authentication:} Reporters can create an account and securely log in to access the platform, and can recover a forgotten password through an email link.
		
		\item \textbf{Report Incident:} Reporters can report an emergency by providing the incident type, severity, title, description, and location, and by attaching up to five photographs.
		
		\item \textbf{Location Selection:} The current location is detected automatically, and the reporter can change it by typing the name of a place or by moving the pin on the map.
		
		\item \textbf{My Incidents:} Reporters can view and filter the incidents they have reported and see their current status.
		
		\item \textbf{Incident Tracking:} Reporters can view the details of an incident, its status timeline, the assigned responder, and the live location of the responder while the response is in progress.
		
		\item \textbf{Edit, Cancel, and Delete:} Reporters can edit or delete an incident while it is pending, and cancel it while it is pending or verified.
		
		\item \textbf{Notifications and Profile:} Reporters receive notifications about status changes and can manage their profile and password.
		
		\item \textbf{Logout:} Reporters can securely end their sessions after using the application.
	\end{enumerate}
	
	\subsection{Admin Module}
	The Admin Module allows the administrator to verify incidents, dispatch responders, and monitor the system.
	
	\begin{enumerate}
		\item \textbf{Admin Authentication:} The administrator can securely log in using valid credentials to access administrative functions.
		
		\item \textbf{Dashboard:} The administrator can view incident counters, the incidents waiting for verification, and the incidents that need a responder.
		
		\item \textbf{Incident Verification:} The administrator can view the details and photographs of a reported incident and verify or reject it.
		
		\item \textbf{Responder Assignment:} The administrator can assign an available responder of the suitable type to a verified incident, and can reassign or cancel the assignment.
		
		\item \textbf{Live Tracking and Nearby Incidents:} The administrator can follow the position of a responder and view active incidents near a chosen location on a map.
		
		\item \textbf{Analytics and Reports:} The administrator can view analytics by incident type, severity, status, and month, and generate a monthly report.
		
		\item \textbf{User Management:} The administrator can create responder accounts, view and edit users, and disable accounts when necessary.
		
		\item \textbf{Logout:} The administrator can securely end the session after completing management activities.
	\end{enumerate}
	
	\subsection{Responder Module}
	The Responder Module allows emergency responders to receive dispatches and report their progress.
	
	\begin{enumerate}
		\item \textbf{Responder Authentication:} Responders log in using the account created for them by the administrator.
		
		\item \textbf{Dashboard and Assignments:} Responders can view their active assignments, completed assignments, and unread notifications.
		
		\item \textbf{Accept or Reject:} Responders can accept an assignment or reject it, after which the incident returns to the administrator for reassignment.
		
		\item \textbf{Location and Navigation:} Responders can view the incident location and photographs, and open the location in Google Maps to navigate to it.
		
		\item \textbf{Response Tracking:} Responders update the status by starting travel, marking arrival, and resolving the incident with optional notes.
		
		\item \textbf{Live Location Sharing:} While an assignment is active, the responder's location is sent to the server at regular intervals.
		
		\item \textbf{History:} Responders can view the incidents they have completed.
		
		\item \textbf{Logout:} Responders can securely end their sessions after use.
	\end{enumerate}
	
	\newpage
	
	% =====================================================================
	%  CHAPTER 6
	% =====================================================================
	\chapter{REPORT OF PROJECT IMPLEMENTATION}
	
	\section{Implementation Plan}
	
	\subsection{Setup and Environment Preparation}
	
	\begin{itemize}
		\item Set up the development environment using Node.js and Visual Studio Code.
		\item Install the required packages for Express.js, Sequelize, MySQL connectivity, authentication, validation, file upload, and the React frontend.
		\item Configure the MySQL database to store users, incidents, assignments, status history, notifications, and tokens.
		\item Configure environment variables for the database, token secrets, upload directory, email service, and the frontend address.
		\item Create the database tables and the initial administrator account using the seed script.
		\item Organize the backend into routes, controllers, services, repositories, models, validators, and middleware.
		\item Organize the frontend into pages, components, API modules, hooks, and a central store for the login session.
	\end{itemize}
	
	\subsection{Backend Implementation}
	
	\begin{itemize}
		\item Develop the backend as a REST API using Express.js, with separate modules for authentication, incidents, responders, administration, users, notifications, and files.
		\item Implement registration, login, logout, token refresh, and password recovery using JSON Web Tokens and hashed refresh tokens.
		\item Implement incident creation, image upload, editing, cancellation, deletion, search, and the geospatial search of nearby incidents.
		\item Implement the incident status workflow and record every status change in the incident history.
		\item Implement responder assignment, reassignment, and cancellation using database transactions with row locking.
		\item Implement responder actions for accepting, rejecting, starting travel, arriving, and resolving, and the endpoint for updating the responder's location.
		\item Generate notifications for important events such as verification, assignment, acceptance, travel, arrival, and resolution.
		\item Implement administrator functions for dashboards, analytics, monthly reports, and user management.
		\item Implement validation of every request using schemas, and a centralized error handler that returns a standard response format.
	\end{itemize}
	
	\subsection{Frontend Implementation}
	
	\begin{itemize}
		\item Develop the user interface using React and Tailwind CSS with a soft and simple design.
		\item Design pages for login, registration, incident reporting, my incidents, incident tracking, responder assignments, the administrator dashboard, analytics, user management, notifications, and profile.
		\item Create the incident reporting form with a map-based location picker, place search, and photograph upload.
		\item Provide the responder with a button that opens the incident location in Google Maps, and buttons for each step of the response.
		\item Display incident and responder locations on maps using Leaflet and OpenStreetMap.
		\item Display analytics charts for the administrator using Recharts.
		\item Fetch and refresh data using TanStack Query, with polling at suitable intervals for notifications and incident status.
		\item Use route guards so that every page is available only to the roles allowed to use it.
		\item Use Tailwind CSS to provide a responsive interface for desktops, laptops, tablets, and smartphones.
	\end{itemize}
	
	\subsection{Security and Data Handling}
	
	\begin{itemize}
		\item Use short-lived access tokens and rotating refresh tokens, and revoke older access tokens when a user's account data changes.
		\item Store passwords as bcrypt hashes and refresh tokens as SHA-256 hashes instead of plain text.
		\item Enforce role-based authorization on every protected route.
		\item Validate request data using schemas, and reject unknown fields.
		\item Use parameterized queries to reduce the risk of SQL injection.
		\item Check the type, size, and file signature of uploaded images, and serve them only to the administrator, the reporter, or the assigned responder.
		\item Hide the contact details of users from roles that do not need them.
		\item Apply security headers, cross-origin restrictions, and rate limiting, with stricter limits on login and password endpoints.
		\item Keep the access token only in the memory of the browser, and keep the refresh token in the browser storage.
	\end{itemize}
	
	\subsection{Project Schedule}
	
	The project was planned over sixteen weeks from July 2026 to October 2026. Table~\ref{tab:schedule} lists the activities and their durations. The key milestones were: requirements completed (end of July), backend completed (end of August), frontend and database integrated (end of September), and testing completed (mid October).
	
	\begin{table}[h!]
		\centering
		\resizebox{\textwidth}{!}{
			\begin{tabular}{|c|p{7cm}|c|c|}
				\hline
				\textbf{ID} & \textbf{Task / Activity} & \textbf{Duration (Weeks)} & \textbf{Weeks} \\
				\hline \hline
				1 & Requirement Analysis & 2 & W1 -- W2 \\
				\hline
				2 & System Design \& Architecture & 2 & W3 -- W4 \\
				\hline
				3 & Backend Development & 4 & W4 -- W7 \\
				\hline
				4 & Frontend Development & 4 & W7 -- W10 \\
				\hline
				5 & Database Design \& Integration & 3 & W9 -- W11 \\
				\hline
				6 & API Integration \& Testing & 2 & W11 -- W12 \\
				\hline
				7 & System Testing \& Bug Fixing & 2 & W13 -- W14 \\
				\hline
				8 & Deployment & 1 & W15 \\
				\hline
				9 & Documentation & 1 & W15 \\
				\hline
				10 & Final Presentation & 1 & W16 \\
				\hline
		\end{tabular}}
		\caption{Project Schedule}
		\label{tab:schedule}
	\end{table}
	
	\section{Testing and Various Types of Testing Used}
	
	Testing verifies that the software functions correctly according to its requirements, helping identify errors and ensure reliable interaction between the frontend, the REST API, and the database. In SafeRoute, testing covers registration and login, role-based access, incident reporting and image upload, verification, responder assignment, the response workflow, location updates, notifications, dashboards, input validation, and access control.
	
	\subsection{Unit Testing}
	
	Unit testing checks individual functionalities independently to verify that each performs its intended task. In SafeRoute, individual features such as registration, login, incident reporting, image upload, verification, assignment, and location update are tested separately. This helps identify errors in individual components before testing the complete application.
	
	\begin{table}[h!]
		\centering
		\resizebox{\textwidth}{!}{
			\begin{tabular}{|c|p{6.5cm}|p{6.5cm}|c|}
				\hline
				\textbf{Test Case} & \textbf{Description} & \textbf{Expected Result} & \textbf{Pass/Fail} \\
				\hline \hline
				1 & Reporter Registration & A new reporter account is created with valid details and a strong password. & Pass \\
				\hline
				2 & User Login & A registered user logs in using valid credentials and receives access and refresh tokens. & Pass \\
				\hline
				3 & Invalid Login & Incorrect login credentials are rejected with an error message. & Pass \\
				\hline
				4 & Report Incident & An incident is saved as pending with its type, severity, and location. & Pass \\
				\hline
				5 & Image Upload & Valid JPEG, PNG, or WebP images are attached, and invalid files are rejected. & Pass \\
				\hline
				6 & Verify Incident & An administrator changes a pending incident to verified. & Pass \\
				\hline
				7 & Assign Responder & An administrator assigns an active responder to a verified incident. & Pass \\
				\hline
				8 & Update Responder Location & The latitude and longitude of the responder are stored. & Pass \\
				\hline
		\end{tabular}}
		\caption{Unit Test Cases and Results}
	\end{table}
	
	\newpage
	
	\subsection{Integration Testing}
	
	Integration testing verifies that different components of SafeRoute work together correctly. It checks the interaction between the frontend, the REST API, and the MySQL database. The testing focuses on whether user actions are processed correctly, whether database changes are reflected in the interface, and whether notifications are generated as expected.
	
	\begin{table}[h!]
		\centering
		\resizebox{\textwidth}{!}{
			\begin{tabular}{|c|p{7cm}|p{7cm}|c|}
				\hline
				\textbf{Test Case} & \textbf{Description} & \textbf{Expected Result} & \textbf{Pass/Fail} \\
				\hline \hline
				1 & Registration and Database Integration & Valid reporter details are stored in the MySQL database. & Pass \\
				\hline
				2 & Login and Token Refresh Integration & The backend verifies credentials, issues tokens, and renews the access token using the refresh token. & Pass \\
				\hline
				3 & Incident Reporting Integration & Incident details, images, and the first status history entry are saved and displayed correctly. & Pass \\
				\hline
				4 & Assignment and Notification Integration & Assigning a responder updates the incident status and notifies the responder and the reporter. & Pass \\
				\hline
				5 & Response Workflow Integration & Accepting, travelling, arriving, and resolving update the incident, the assignment, and the history. & Pass \\
				\hline
				6 & Secured Image Access & Photographs are displayed only to the administrator, the reporter, and the assigned responder. & Pass \\
				\hline
				7 & Dashboard and Analytics Integration & Counters, charts, and reports show data retrieved from the database. & Pass \\
				\hline
				8 & Notification Integration & Notifications can be listed, marked as read, and deleted, and the unread count is updated. & Pass \\
				\hline
		\end{tabular}}
		\caption{Integration Test Cases and Results}
	\end{table}
	
	\newpage
	
	\subsection{System Testing}
	
	System testing evaluates the complete SafeRoute application to verify that its major features operate together as expected. It includes testing the complete flow from reporting an incident to its resolution with the three user roles, session handling, access restrictions, location features, and the responsiveness of the interface. It also checks input validation and error handling to identify problems that may affect the overall reliability and usability of the application.
	
	\begin{table}[h!]
		\centering
		\resizebox{\textwidth}{!}{
			\begin{tabular}{|c|p{6cm}|p{6cm}|c|}
				\hline
				\textbf{Test Case} & \textbf{Description} & \textbf{Expected Result} & \textbf{Pass/Fail} \\
				\hline \hline
				1 & Complete Incident Workflow & An incident moves from reporting to verification, assignment, acceptance, travel, arrival, and resolution. & Pass \\
				\hline
				2 & Authentication and Session Testing & Valid credentials are accepted, invalid credentials are rejected, and expired sessions are renewed. & Pass \\
				\hline
				3 & Role-Based Access Testing & Users cannot access pages or operations that belong to another role. & Pass \\
				\hline
				4 & Reporter Incident Management & Reporters can edit, cancel, and delete incidents only in the allowed statuses. & Pass \\
				\hline
				5 & Location and Navigation Testing & Incident locations are displayed on the map, and the Google Maps button opens the correct destination. & Pass \\
				\hline
				6 & Admin Management Testing & The administrator can verify, assign, reassign, and manage users and view analytics. & Pass \\
				\hline
				7 & Input Validation Testing & Invalid or incomplete input is rejected with appropriate messages. & Pass \\
				\hline
				8 & Responsiveness Testing & The interface works correctly on desktop and mobile screen sizes. & Pass \\
				\hline
		\end{tabular}}
		\caption{System Test Cases and Results}
	\end{table}
	
	% =====================================================================
	%  CHAPTER 7
	% =====================================================================
	\chapter{RESULTS AND DISCUSSION}
	
	\section{Advantages}
	
	\begin{itemize}
		\item \textbf{Centralized Platform:}
		SafeRoute connects reporters, administrators, and responders through a single system for reporting, verification, dispatch, and resolution.
		
		\item \textbf{Faster and More Detailed Reporting:}
		Reporters can submit an incident with its type, severity, location, and photographs in one form.
		
		\item \textbf{Verified Incidents:}
		Administrators verify incidents before dispatch, which reduces the effort spent on false or duplicate reports.
		
		\item \textbf{Organized Dispatch:}
		Administrators can assign a suitable responder, and reassign or cancel the assignment when required.
		
		\item \textbf{Navigation Support:}
		Responders can open the incident location directly in Google Maps, which helps them reach the location.
		
		\item \textbf{Real-Time Status Tracking:}
		Reporters and administrators can follow the status of an incident and the position of the responder during the response.
		
		\item \textbf{Notifications:}
		Users are informed about assignments and status changes without having to ask for updates.
		
		\item \textbf{Complete Incident History:}
		Every status change is recorded with the user and the time, which supports monitoring, transparency, and analysis.
		
		\item \textbf{Role-Based Security:}
		Each user can access only the functions meant for their role, and incident photographs and contact details are protected.
		
		\item \textbf{Free and Open Mapping:}
		The use of Leaflet and OpenStreetMap avoids the need for paid map services.
	\end{itemize}
	
	\newpage
	
	\section{Limitations}
	
	\begin{itemize}
		\item \textbf{Internet Dependency:}
		The web application requires an internet or network connection to report incidents, receive assignments, and use location services.
		
		\item \textbf{Polling Instead of Push:}
		Notifications and live status are refreshed at regular intervals, so updates may appear after a short delay instead of instantly.
		
		\item \textbf{Live Location Depends on the Device:}
		The responder's location is shared only while the page remains open and location access is allowed on the device.
		
		\item \textbf{Manual Dispatch Decisions:}
		Verification and assignment are performed by the administrator, so the response time depends on how quickly the administrator acts.
		
		\item \textbf{Dependence on Reporter Information:}
		The usefulness of a report depends on the accuracy of the location and details given by the reporter.
		
		\item \textbf{Address Search Service Limits:}
		Place search and address lookup use a free public service that has usage limits.
		
		\item \textbf{Local Image Storage:}
		Photographs are stored on the server's disk, so storage and backups must be managed as the number of incidents grows.
		
		\item \textbf{Scalability Challenges:}
		Supporting a very large number of users and incidents may require further performance optimization and infrastructure improvements.
	\end{itemize}
	
	\section{Screenshots}
	\vspace{0.1cm}
	
	\begin{figure}[H]
		\centering
		\includegraphics[width=0.8\linewidth,height=0.6\textheight,keepaspectratio]{Screenshot-login.png}
		\caption{Login Page}
	\end{figure}
	
	\begin{figure}[H]
		\centering
		\includegraphics[width=0.8\linewidth,height=0.6\textheight,keepaspectratio]{Screenshot-reporter-registration.png}
		\caption{Reporter Registration}
	\end{figure}
	
	\begin{figure}[H]
		\centering
		\includegraphics[width=0.8\linewidth,height=0.6\textheight,keepaspectratio]{Screenshot-my-incidents.png}
		\caption{Reporter -- My Incidents}
	\end{figure}
	
	\begin{figure}[H]
		\centering
		\includegraphics[width=0.8\linewidth,height=0.6\textheight,keepaspectratio]{Screenshot-report-emergency.png}
		\caption{Reporter -- Report Emergency}
	\end{figure}
	
	\begin{figure}[H]
		\centering
		\includegraphics[width=0.8\linewidth,height=0.6\textheight,keepaspectratio]{Screenshot-responder-dashboard.png}
		\caption{Responder Dashboard}
	\end{figure}
	
	\begin{figure}[H]
		\centering
		\includegraphics[width=0.8\linewidth,height=0.6\textheight,keepaspectratio]{Screenshot-responder-emergency-view.png}
		\caption{Responder -- Emergency Details and Response Actions}
	\end{figure}
	
	\begin{figure}[H]
		\centering
		\includegraphics[width=0.8\linewidth,height=0.6\textheight,keepaspectratio]{Screenshot-responder-assignment-history.png}
		\caption{Responder -- Assignment History}
	\end{figure}
	
	\begin{figure}[H]
		\centering
		\includegraphics[width=0.8\linewidth,height=0.6\textheight,keepaspectratio]{Screenshot-admin-dashboard.png}
		\caption{Admin Dashboard}
	\end{figure}
	
	\begin{figure}[H]
		\centering
		\includegraphics[width=0.8\linewidth,height=0.6\textheight,keepaspectratio]{Screenshot-admin-live-map-view.png}
		\caption{Admin -- Live Map}
	\end{figure}
	
	\begin{figure}[H]
		\centering
		\includegraphics[width=0.8\linewidth,height=0.6\textheight,keepaspectratio]{Screenshot-admin-analytics-view.png}
		\caption{Admin -- Analytics}
	\end{figure}
	
	\begin{figure}[H]
		\centering
		\includegraphics[width=0.8\linewidth,height=0.6\textheight,keepaspectratio]{Screenshot-admin-responder-create-view.png}
		\caption{Admin -- Create Responder Account}
	\end{figure}
	
	\newpage
	
	% =====================================================================
	%  CHAPTER 8
	% =====================================================================
	\chapter{CONCLUSION AND FUTURE SCOPE}
	
	\section{Conclusion}
	
	\hspace{1cm}The development of SafeRoute provides a centralized approach to emergency incident reporting and response management. By connecting reporters, administrators, and responders through a single system, it enables faster incident reporting and better coordination during emergencies. Incident verification, responder assignment, and status tracking improve the organization of the response, and the centralized incident records provide better monitoring, transparency, and response history.
	
	The platform includes features such as role-based registration and login, incident reporting with location and photographs, administrator verification and dispatch, responder acceptance, navigation, and live location sharing, notifications, and dashboards and reports for administrators. The proposed system can help reduce communication gaps and delays in emergency response, and it provides a foundation for future improvements such as advanced location services, automated alerts, and intelligent response features.
	
	\section{Future Scope}
	
	\begin{enumerate}
		\item \textbf{Mobile Application Development:}
		Developing a mobile application so that reporters and responders can use the platform conveniently from their smartphones, including reliable background location sharing.
		
		\item \textbf{Real-Time Push Updates:}
		Replacing periodic polling with real-time updates and push notifications so that assignments and status changes reach users instantly.
		
		\item \textbf{Automated Alerts:}
		Introducing automated SMS and email alerts for critical incidents and for assignments that are not accepted within a given time.
		
		\item \textbf{Intelligent Responder Suggestion:}
		Suggesting the most suitable responder automatically, based on the type of the incident and the distance from the responder to the location.
		
		\item \textbf{Advanced Location Services:}
		Adding in-application routing, estimated arrival times, and live tracking of responders along their route.
		
		\item \textbf{Offline Reporting:}
		Allowing incidents to be saved on the device when there is no network and sent automatically when the connection is restored.
		
		\item \textbf{Advanced Analytics:}
		Adding heat maps, response-time analysis, and exportable reports to help authorities plan and improve their response.
	\end{enumerate}
	
	% (Add a code appendix here later if needed.)
	
	\vspace{2cm}
	
	\renewcommand{\bibname}{\uppercase{REFERENCES}}
	
	\begin{thebibliography}{999}
		
		\addcontentsline{toc}{chapter}{\hspace{0.19in} REFERENCES}
		
		\bibitem{one}
		S. Zhou, M. Wu, X. Ye, J. Silbernagel, D. A. Hart, and Q. Huang, ``Advancing smart disaster response by leveraging social sensing and mobile technology,'' \textit{Geo-Spatial Information Science}, vol. 29, no. 4, pp. 2651--2669, 2026.
		
		\bibitem{two}
		W. Alkusaibati, S. Pilemalm, T. Andersson Granberg, E. Prytz, and N. Ennab Vogel, ``From dispatch to impact: Evaluating ICT-enabled co-production of emergency response,'' \textit{International Journal of Disaster Risk Reduction}, vol. 130, Art. no. 105833, 2025.
		
		\bibitem{three}
		S. K. Sharma, J. B. Singh, and M. Kumar, ```Invented-on-the-fly' mobile application for disaster response: Construction of technological frames and impact,'' \textit{Information Systems Journal}, vol. 35, no. 4, pp. 1132--1165, 2025.
		
		\bibitem{four}
		M. Syukron, A. Madugalla, S. Mojtaba, and J. Grundy, ``Engineering for crisis management: A user-centred analysis of disaster mobile applications,'' \textit{SSRN}, 2026.
		
	\end{thebibliography}
	
\end{document}